%%=====================================================================
%% M2 — WP5: what ω is. SURGICAL FRAGMENT. UK English. "We".
%% PI rulings: Q4 (revised; the two zeros derived in M2, GUE maximum [M],
%% spectral proof to M3/M4), Q3 (half-circle route (a) after R3.1), R1
%% (ω = ψ + ψ̄). Depends on WP0 (def:omega, prop:parity), WP2 (thm:handedness,
%% M1 Thm 4.2 drive), WP4 (eq:a-Phi, decaying warp). Placeholder numbering.
%%=====================================================================

\section{What the Coherent Field is}\label{sec:omega-shape}

The Coherent Field entered in Definition~\ref{def:omega} as the even
combination of the sheets. This section says what it measures. The
programme reads $\omega$ as a one-hump quantity: zero when the sheets are
rigidly locked, zero when they are uncorrelated, positive between. We
derive the two zeros from A2 and A4, show where the maximum of the
coupling sits, settle the sign of $\omega$ on the solutions of
Section~\ref{sec:warp}, and name what this paper does not prove.

\subsection{The phase of a sheet field}\label{sec:phase}

Work in the corner reading of Definition~\ref{def:psi}, where each
component is complex, and write $\psi=re^{i\vartheta}$ in $\mathbb{C}_{u}$,
so that $\omega=2r\cos\vartheta$ and $\psi-\bar\psi=2ir\sin\vartheta$; the
statements below hold componentwise, and the module trace sums them. The phase
$\vartheta$ is the angle between the two sheets: $\vartheta=0$ when they
coincide, $\vartheta=\pm\pi/2$ when growth stands exactly against decay.
By Definition~\ref{def:psi}, $\hat\tau$ sends $\vartheta\mapsto-\vartheta$.

\begin{proposition}[the locked zero]\label{prop:locked}
$\omega=0$ at a point iff $\bar\psi=-\psi$ there, i.e. iff the sheet field
is $\hat\tau$-odd, $\vartheta=\pm\pi/2$. At such a point the dilaton
vanishes, $\QCD=0$, and the geometry does not respond to the sheets.
\end{proposition}

\begin{proof}
The Coherent Field is $\psi$ plus its shadow. It is zero at a point exactly
when the shadow there is the negative of the field, $\hat\tau\psi=-\psi$:
when exchanging growth and decay reverses the sign of every component. In
the corner reading, where a component is $re^{i\vartheta}$ and the exchange
sends $\vartheta$ to $-\vartheta$, this says $re^{i\vartheta}=-re^{-i\vartheta}$,
that is $\cos\vartheta=0$, so $\vartheta=\pm\pi/2$: the two sheets stand at
right angles, growth exactly against decay. At such a point the dilaton
$\QCD=\omega\lvert\psi-\bar\psi\rvert^{2}$ is zero because its first factor is,
and so the $\zeta$-term of \eqref{eq:master} vanishes there: the geometry
does not respond. \rigR
\end{proof}

This is A2 read on the field: when growth is exactly opposed by decay the
two sheets are locked, nothing in the geometry changes, and the even
combination is zero. It is the reason the Coherent Field is the
combination $\psi+\bar\psi$ and not the modulus, which does not vanish
there (Section~\ref{sec:omega}). We call this the \emph{locked} or rigid
configuration.
%% PI ruling R1, now a proposition.

\subsection{The uncorrelated zero, from the drive}\label{sec:uncorrelated}

The second zero comes from A4. The drive $M_{0}$ acts on the base torus
(M1 Theorem~4.2), and a sheet field is carried along its orbit,
$\psi\mapsto\psi\circ M_{0}^{n}$. We make precise the average that the
programme's ``orbit average'' refers to.

\begin{definition}[orbit average]\label{def:orbit-average}
For a sheet field $\psi$ and a point $z$ of the base, the orbit average of
$\omega$ over $T$ circuits is
$\langle\omega\rangle_{T}(z)=\tfrac1T\sum_{n<T}\omega(M_{0}^{n}z)$. The
\emph{phase distribution} of $\psi$ is the push-forward of Lebesgue
measure on the base under $z\mapsto\vartheta(z)$.
\end{definition}

\begin{proposition}[the uncorrelated zero]\label{prop:uncorrelated}
The drive is ergodic for Lebesgue measure, so for almost every $z$,
$\langle\omega\rangle_{T}(z)\to\int\omega\,\mathrm{d}\mu$ as $T\to\infty$:
the orbit average is the space average. If the phase distribution of
$\psi$ is uniform on the circle --- the phase of the sheets is
uncorrelated with position --- then $\int\omega\,\mathrm{d}\mu=0$.
\end{proposition}

\begin{proof}
A hyperbolic toral automorphism is ergodic, indeed mixing, for Lebesgue
measure \cite{KatokHasselblatt}; M1 Proposition~4.1 shows A4 forces
hyperbolicity. Birkhoff's theorem gives the first statement. For the
second, with $r$ bounded and $\vartheta$ uniform,
$\int\omega=2\int r\cos\vartheta\,\mathrm{d}\mu=2\bar r\int_{0}^{2\pi}\cos\vartheta\,\tfrac{\mathrm{d}\vartheta}{2\pi}=0$,
where the factorisation uses that a uniform phase is independent of the
modulus. \rigR\ (given the independence, which is what ``uniform phase''
means).
\end{proof}

This is A4 read on the field: a drive of positive entropy mixes, and a
sheet field whose phase has been mixed to uniformity has no net even
combination. We call this the \emph{uncorrelated} configuration.

\subsection{The hump, and where the coupling is largest}\label{sec:hump}

Between the two zeros the average is positive. Let the phase distribution
be concentrated around a lock angle $\vartheta_{0}$ with concentration
$c\ge0$ (any unimodal family will do; the circular normal
$\propto e^{c\cos(\vartheta-\vartheta_{0})}$ is the standard one). Then
\begin{equation}\label{eq:hump}
  \langle\omega\rangle\;=\;2\bar r\,A(c)\cos\vartheta_{0}\,,\qquad
  A(c)=\frac{I_{1}(c)}{I_{0}(c)}\in[0,1)\,,
\end{equation}
which vanishes at $c=0$ (uniform, Proposition~\ref{prop:uncorrelated}),
vanishes at $\vartheta_{0}=\pm\pi/2$ (locked,
Proposition~\ref{prop:locked}), and is positive for
$\lvert\vartheta_{0}\rvert<\pi/2$: one hump in the lock angle between the
two locked configurations. \rigD

The coupling to curvature is not maximal where $\omega$ is. In the phase,
\begin{equation}\label{eq:Phi-phase}
  \QCD\;=\;\omega\,\lvert\psi-\bar\psi\rvert^{2}\;=\;8r^{3}\cos\vartheta\,\sin^{2}\vartheta\,,
\end{equation}
which vanishes at $\vartheta=0$ (sheets coincident, no held distinction)
and at $\vartheta=\pm\pi/2$ (locked), and is extremal where
$\tan^{2}\vartheta=2$, $\sin^{2}\vartheta=\tfrac23$. The geometry responds
most strongly neither when the sheets agree nor when they oppose, but at
the definite intermediate angle $\vartheta^{\ast}=\arctan\sqrt2\approx54.7^{\circ}$
(\texttt{L6\_Phi\_extremum}). \rigR
%% Direct computation from prop:parity; the angle is the parity theorem's
%% own number. Not in the corpus as far as read; flagged as new.

\subsection{The sign of $\omega$ on the solutions}\label{sec:half-circle}

Section~\ref{sec:warp} gave the warp as a decaying exponential,
$a=-\zeta\QCD$ with $\QCD$ rising from $0$ at the wall to $1$ at the far
wall, $\QCD$ there standing for the real part of the trace. On those
solutions, therefore, $\QCD\ge0$. But $\QCD$ is $\omega$ times a
non-negative real number, $\lvert\psi-\bar\psi\rvert^{2}$, so wherever the
sheets are distinct --- wherever that number is not zero --- the real part
of $\omega$ has the same sign as $\QCD$ and is non-negative: in the corner
reading, $\cos\vartheta\ge0$, which is the half-circle
$\lvert\vartheta\rvert\le\pi/2$. The Coherent Field does not change sign on
the solutions. What fixed the sign was the direction of the warp: the
exponential decays away from the wall, so $\QCD$ rises. Had the warp grown
instead, the half-circle would be the other one, and the overall sign of
$\psi$ --- which no axiom fixes, since $\psi\mapsto-\psi$ commutes with
$\hat\tau$ and reverses $\omega$ and $\QCD$ together --- would be chosen to
restore $\mathrm{Re}\,\omega\ge0$. Away from the corner reading $\omega$
carries an imaginary part along the selected unit, which the sign statement
does not constrain. \rigD
%% Closes R2.4 (Q3, route (a)). PI ruling 9 Nov 2026: modulus squared; the
%% signs of Φ and Re ω now agree directly. Minimality Register M2-9 updated.

\subsection{What this paper does not prove}\label{sec:omega-hook}

The programme's canonical reading of $\omega$ is a one-hump complexity
functional of a \emph{spectrum}: zero for a rigidly ordered spectrum,
zero for an uncorrelated one, maximal in the regime of strong level
repulsion, with spectral rigidity as the computable proxy. The two zeros
derived above are the field-level counterparts of the first two
statements: locked for rigidly ordered, uniform for uncorrelated. The
identification of the field $\omega$ with that spectral functional, and
the placement of its maximum in the regime of strong level repulsion, are
not proved here; they require the spectrum of the Dirac operator at the
cone points of the pillowcase and a bridge theorem between the two
descriptions, which is the subject of a later paper of the series. We
record the reading as motivated and the bridge as open. \rigM
%% Q4 (revised): GUE maximum [M]; spectral proof registered as M3/M4
%% (Register §12.1). ω named per omega-definition-usage §6a throughout.
